\section{Nonclassical Light}
All light states are nonclassical, but it can be hard to see the quantum mechanical effects. The only pure light state that have a $P(\alpha)$ distribution that is positive everywhere and note more singular than a delta function is the coherent state.

\subsection{Quadrature Squeezing}
Two operators $\hat{A},\hat{B}$ with commutation relation $[\hat{A},\hat{B}] = i \hat{C}$ will have $\langle (\Delta \hat{A})^2 \rangle\langle (\Delta \hat{B})^2 \rangle \geq |\langle\hat{C}\rangle|^2 /4$. We call the system squeezed if $\langle (\Delta \hat{A})^2 \rangle < |\langle\hat{C}\rangle|/2$ (or similar for $\hat{B}$). Ideal squeezed states are states for which $[\X, \XX] = i/2$ holds. Quadrature squeezing exist when $\hat{A} = \X, \hat{B} = \XX, \hat{C} = 1/2$ and $\langle (\Delta\X)^2\rangle < 1/4$ (or similar for $\XX$). For this to hold we require that the $P$-function is nonpositive somewhere. We can introduce a general quadrature operator $\Xgen(\theta) = [\a\exp(i\theta) + \ad\exp(-i\theta)]/2$ from where $\X$ and $\XX$ follow. For the parameter $s(\theta) = 4\langle(\Delta \Xgen(\theta))^2\rangle - 1$, we have squeezed states when $-1 \leq s(\theta) < 0$.

We can generate squeezed states by introducing a squeezing operator $\S(\xi) = \exp[(\xi^*\a^2 - \xi (\ad)^2)/2]$, where $\xi = r\exp(i\theta)$, with $r$ the squeeze parameter, and $0 \leq r < \infty, 0\leq\theta\leq 2\pi$. \question{Where does this operator come from?} The squeezing operator can be interpreted as a two-photon displacement operator. We clearly have $\Sd(\xi) = \S(-\xi)$ and the Baker-Housdorff lemma gives $\Sd(\xi)\a\S(\xi) = \a\cosh r - \ad \exp(i\theta) \sinh r$ and $\Sd(\xi)\ad\S(\xi) = \ad\cosh r - \a \exp(-i\theta) \sinh r$. Applying the squeezing operator to the vacuum state we get $\ket{\xi} = \S(\xi)\ket{0}$, and calculating the fluctuations of the quadrature operator for $\theta = 0$ yields $\lineexpval{(\Delta \X)^2} = \exp(-2r)/4$ and $\lineexpval{(\Delta \XX)^2} = \exp(2r)/4$. Thus, we have squeezing in the $\X$ quadrature, choosing $\theta = \pi$ we get squeezing in the $\XX$ quadrature, for an arbitrary $\theta$ we can rotate the coordinate system and find a similar result. Applying also the displacement operator $\ket{\alpha, \xi} = \D(\alpha)\S(\xi)\ket{0}$ we obtain a squeezed state that is displaced from the origin by $\alpha$. This can be expanded as $\ket{\alpha, \xi} = \sum_n c_n \ket{n}$ with 
\begin{equation}
    c_n = \frac{1}{\sqrt{\cosh(r)}}\exp[-\frac{1}{2}\abs{\gamma}^2 + \frac{1}{2}\gamma^2 e^{-i\theta}\tanh(r)] \frac{1}{\sqrt{n!}} \left[\frac{1}{2} e^{i\theta}\tanh(r)\right]^{n/2} H_n \left[\gamma /\sqrt{e^{i\theta}\sinh(2r)}\right],\label{eq:cnSqueezedDisplaced}
\end{equation}
with $\gamma = \alpha\cosh(r) + \alpha^*e^{i\theta}\sinh(r)$ and $H_n$ the Hermite polynomials.

The electric field of a squeezed state will have phase-dependant noise. Similar to how we generated an electric field by rotating the coherent state around the origin, with the uncertainty in the field coming from the uncertainty in the coherent state's quadratures. Since a squeezed state has less uncertainty in one of the quadratures and more in the other, the projection while rotating around the origin will have very different uncertainties. For example, this can be used to obtain an electric field where the peak position has very little uncertainty. This fact can be used in the detection of very weak signals, like gravity wave detector. \study{It would be interesting to use these state in quantum sensing applications.}

A squeezed state on the form $\ket{\alpha, \xi}$ can be characterized by the correlation between the quadratures. Also, important to note is that the squeezed state $\ket{\alpha, \xi}$ is not every possible squeezed state. \question{What other types of squeezed states are there?}

\subsection{Generation of Quadrature Squeezed Light}
To generate quadrature squeezed light most schemes use degenerate parametric down-conversion, realized by nonlinear optics. If a certain nonlinear medium is pumped by a laser with frequency $\omega_p$ it produces two photons with frequency $\omega / \omega_p/2$ with the Hamiltonian $\hamiltonian = \hbar\omega\ad\a + \hbar\omega_p\hat{b}^\dagger \hat{b} + i\hbar \chi^{(2)} (\a^2 \hat{b}^\dagger - (\ad)^2 \hat{b})$, with $a$ the signal mode, $b$ the pump mode, and $\chi^{(2)}$ is the second order susceptibility. The parametric approximation says that the pump field is in a strong classical coherent field that remains undepleted of photons during the relevant timescale. Assuming the coherent state and approximation $\hat{b} = \beta \exp(-i\omega_p t)$ we can write the Hamiltonian in the interaction picture $\hamiltonian_I (t) = i\hbar\{ \eta^* \a^2 \exp[i(\omega_p - 2\omega)t] - \eta (\ad)^2 \exp[-i(\omega_p - 2\omega)t] \}$ with $\eta = \beta \chi^{(2)}$. The time dependency may be removed by choosing a material such that $\omega_p = 2\omega$ giving the Hamiltonian $\hamiltonian_I = i\hbar (\eta^* \a^2 - \eta (\ad)^2)$, with the unitary evolution operator $\hat{U}(t,0) = \exp(\eta^* t \a^2 - \eta t (\ad)^2)$ which is the squeezing operator for $\xi = 2\eta t$. There is a similar process called degenerate four-wave mixing that may also be used.

\subsection{Detection of Quadrature Squeezed Light}
To detect or measure squeezed states we can use a method called balanced homodyne detection. We use a 50/50 beam splitter with the $\a$ input state being the signal we want to measure, and $\hat{b}$ a strong coherent laser field with the same frequency as the signal (homodyne) but with adjustable phase. The two output ($\hat{c},\hat{d}$) intensities of the beam splitter is measured, and a correlator takes the difference. Assuming the coherent state $\ket{\beta e^{-i\omega t}}$ and $\beta = |\beta|\exp(i\psi)$ with $\psi$ the phase, and defining $\theta = \psi + \pi/2$, we get that the number difference in the detectors is  $\lineexpval{n_{cd}} = \lineexpval{\hat{c}^\dagger \hat{c} - \hat{d}^\dagger\hat{d}} = 2|\beta| \lineexpval{\Xgen(\theta)}$. With this we can measure and arbitrary quadrature by tuning the phase of the laser (changing $\theta$).

\subsection{Amplitude (or Number) Squeezed States}
The heuristic number phase commutator $[\hat{n}, \hat{\phi}] = i$ lead to the uncertainty relation $\Delta n \Delta \phi \geq 1/2$. Thus, by analogy to the quadrature squeezing we can imagine states squeezed in number such that $\Delta n < \sqrt{\bar{n}}$ or in phase such that $\Delta \phi < 1 / (2\sqrt{\bar{n}})$. We note that there does not exist a Hermitian phase operator, which create problems. We can still define some photon statistics for amplitude squeezed states, which are nonclassical and sub-Possonian. \question{I don't like using the phase like this since we lack a phase-operator. It feels weird, and I do not understand why we can do this.} A way to categorize states is with Mandel's $Q$-parameter $Q = (\lineexpval{(\Delta \hat{n})^2 - \lineexpval{\hat{n}}})/\lineexpval{\hat{n}}$. For $-1 \leq Q < 0$ the states are said to be sub-Poissonian and for $Q = 0$ the states are Poissonian, and for $Q > 0$ super-Poissonian. A coherent state have $Q=0$ and thus sub-Poissonian states have a narrower distribution and super-Poissonian states have a broader distribution than a coherent state. It is also possible for a state to be both number squeezed and quadrature squeezed.

\subsection{Photon Antibunching}
Antibunching is when the probability for joint detection increases with time which is a nonclassical effect, that is $\qqcoherence(\tau) > \qqcoherence(0)$. The second order coherence function is related to the $Q$-parameter for a single mode field by $Q = \lineexpval{\hat{n}}[\qqcoherence(0) - 1]$ and since $\qqcoherence(\tau) = \qqcoherence(0)=\mathrm{constant}$ we can not see bunching or antibunching effects for a single mode field, since we lack the time dependence. To see this we need multimode fields which will give the time dependence. For a driving field with Rabi frequency $\Omega$ and spontaneous emission decay rate $\gamma$, for $\Omega \gg \gamma$ we get $\qqcoherence(\tau) = 1 - \exp(-3\gamma\tau/4)\cos(\Omega t)$ and for $\Omega \ll \gamma$ we get $\qqcoherence(\tau) = [1- \exp(-\gamma\tau/2)]^2$. Both of which have $\qqcoherence(0) = 0$ and $\qqcoherence(\tau) > 0 $ for all $\tau > 0$, and thus showing antibunching.

\subsection{Schrödinger Cat-States}
A Schrödinger cat state is a state on the form $\ket{Psi} = N (\ket{\alpha} + \exp(i\Phi)\ket{-\alpha})$ where $N$ is a normalization constant, with large $|\alpha|$ the states are macroscopically distinguishable. The question is why we cannot see macroscopic entanglement in everyday life (some recent experiments have managed to create smaller macroscopic states of this form)? No quantum state is truly isolated, especially not a macroscopic one, and thus, if a pure macroscopic state is created it would quickly interact with the environment and decohere into a classical statistical mixture, or a mixed state. Depending on the chosen phase there are three main types of cat states. For $\Phi = 0$ we have the even states $\ket{\Psi_e} = N_e (\ket{\alpha} + \ket{-\alpha})$, for $\Phi = \pi$ we have the odd states $\ket{\Psi_o} = N_o (\ket{\alpha} - \ket{-\alpha})$, and finally for $\Psi = \pi/2$ we have the Yurke-Stoler states $\ket{\Psi_{YS}} = (\ket{\alpha} + i\ket{-\alpha})/\sqrt{2}$. All of these states are eigenstates to $\a^2$ with eigenvalue $\alpha^2$. 

The density matrix for the statistical mixture is defined as $\density_\text{mixture} = (\ket{\alpha}\bra{\alpha} + \ket{-\alpha}\bra{-\alpha})/2$, and the (for example) even cat state is $\density_e = |N_e|^2 (\ket{\alpha}\bra{\alpha} + \ket{\alpha}\bra{-\alpha} + \ket{-\alpha}\bra{\alpha} + \ket{-\alpha}\bra{-\alpha})$. To distinguish between these two states we must be able to detect the coherences. This may be done using photon statistics as for the even cat state $P_n = 0$ for $n$ odd while $P_n$ is nonzero for $n$ even, while for the mixed state it is nonzero for all $n$, a similar argument can be done for the odd cat state. The Yurke-Stoler state is also nonzero for all $n$. Both the even cat state and the Yurke-Stoler cat state exhibit quadrature squeezing in $\XX$, while there is no squeezing for the odd cat state and the mixed state. \question{How are cat states used as qubits? I've heard that they can be.}

The Lindblad (CKSL) master equation is $\partial_t \density = -i /\hbar \times [\hamiltonian, \density(t)]  + \sum_\alpha \gamma_\alpha (\L_\alpha \density(t) \Ld_\alpha - \{ \Ld\L, \density(t) \}/2)$, where $\L$ is jump operator and $\gamma$ a parameter related to the probability (or rate) of the jump occurring and can either be calculated with Fermi's golden rule or fit to experimental data. The jump operator is an operator that transfer one quantum state into another. Examples can be the loss or gain of a photon in a cavity or the excitation or deexcitation of an atom. \study{I would like to study the derivation of the Lindblad equation further as I still do not feel completely comfortable with it.}

\subsection*{Appendix D: Nonlinear Optics and Spontaneous Parametric Down-Conversion}\addcontentsline{toc}{subsection}{Appendix D: Nonlinear Optics and Spontaneous Parametric Down-Conversion}
The physical effect of down-conversion is due to the fact that in nonlinear optics the induced polarization is strong enough to need at least second order susceptibility to be considered. That is the linear response in polarization is not sufficient for these materials. Noting here that many materials used for down-conversion are anisotropic. \question{Does there exist materials that can be used for down-conversion that are isotropic?} The parametric approximation tells us that we pump with a very strong coherent laser field that can be treated classically as long as the interactions are short enough that the pump photons don't run out. Doing this we get the phase matching conditions $\omega_p = \omega_s + \omega_i$ and $\vec{k}_p = \vec{k}_s + \vec{k}_i$, where $p$ is for pump, $s$ for signal, and $i$ for idler.